\cotefile{63}
\scanpage{25}
\mspage{10}
\makeatletter\tagsleft@true\makeatother

(3.1), $c_4=0$ si et seulement si $a_2^2=0$. Supposons qu'on ait même $a_2=0$
\[
\begin{aligned}[t]
 &y^2=x^3+a_4 x+a_6 \qquad \pi=a_4^{-1}\,dy=dx/2y\\[.65em]
 &c_4=c_6=0 \qquad \Delta=-a_4^3 \, .
\end{aligned}\tag*{(5.1)}
\]

Les autres systèmes $(\pi,x,y)$ donnant une équation de cette forme s'obtiennent en faisant $s=t=0$ dans (1.2) et
\[
\begin{aligned}
u^4 {\unsure{\mathbb{\text{\makebox[0pt][l]{\kern.45pt$a$}$a$}}|a}}'_4 &= a_4\\
u^6 a'_6 &= a_6+r\,a_4+r^3 \, .
\end{aligned}\tag*{(5.2)}
\]

b) $p=2$\tikz[baseline=(p.base)]{\node[inner sep=0pt] (p) {.};\handwritten{\draw[thin,gray] (0.13em,1.1em) .. controls (-0.42em,1.25em) and (-0.57em,-1.52em) .. (-0.01em,-1.45em) .. controls (0.40em,-1.46em) and (0.77em,0.33em) .. (0.75em,1.44em);}}

En caractéristique 2 , on a $c_4=a_1^4$. Supposons qu'on ait $a_1=0$ ; $\pi$ étant donnée, on peut alors choisir $x$ et $y$ tels que (1.3) soit
\[
\begin{aligned}[t]
 &y^2+a_3\,y=x^3+a_4\,x+a_6 \qquad
\pi=\frac{dx}{a_3}=dy/\tikz[baseline=(den.base)]{\node[inner sep=0pt] (den) {$x^2+a_4$};\handwritten{\draw[<->,thin] ([yshift=0.6ex]den.north west) .. controls ([xshift=-0.3em,yshift=-1.1ex]den.south west) and ([xshift=-0.1em,yshift=-1.1ex]den.south east) .. ([xshift=0.32em,yshift=-0.35ex]den.south east) .. controls ([xshift=0.6em]den.south east) and ([xshift=0.32em]den.north east) .. ([xshift=0.32em,yshift=0.65ex]den.north east);}}
\\[.65em]
 &b_2=b_4=0 \qquad b_6=a_3^2 \qquad b_8=a_4^2 \qquad
\Delta=a_3^4 \qquad c_4=c_6=0 \, .
\end{aligned}\tag*{(5.3)}
\]

Les autres systèmes $(\tikz[baseline=(pc.base)]{\node[inner sep=0pt] (pc) {$\pi,$};\handwritten{\draw[thin,gray,rotate around={-9:(pc.center)}] (pc.center) ellipse (0.62em and 1.07em);}}\,x,y)$ donnant une équation de cette forme s'obtiennent en faisant $r=s^2$ dans (1.2) et
\[
\begin{aligned}
u^3 a'_3 &= a_3\\
u^4 a'_4 &= a_4+sa_3+s^4\\
u^6 a'_6 &= a_6+s^2\,a_4+ta_3+s^6+t^2 \, .
\end{aligned}\tag*{(5.4)}
\]
% archive sheet number: 24

\scanpage{26}
\mspage{11}

6. \quad \uline{Applications}.

\uline{Proposition} 6.1. \quad \uline{Si} $E$ \uline{est une courbe de genre un sur un corps algébriquement clos}, $E$ \uline{est de type elliptique} (\uline{resp. multiplicatif}, \uline{resp. additif}) \uline{si et seulement si} $\Delta\ne 0$ (\uline{resp.} $\Delta=0$ et $c_4\ne 0$, \uline{resp}. $\Delta=c_4=0$).

a) \quad Lorsque la caractéristique $p\ne 2$, la première assertion résulte de la troisième formule (3.2). Lorsque $p=2$ et $c_4\ne 0$, on vérifie aussitôt que la différentielle de (4.1) ne s'annule qu'en $(0,0)$, et $\Delta$ est elliptique si et seulement si $\Delta=a_6\ne 0$. Pour $p=2$ et $c_4=0$, on raisonne de même sur (6.2).

b) \quad Lorsque $p\ne 2$, il résulte de (3.2) que $\Delta=c_4=0$ si et seulement si le deuxième membre de (3.1) a toutes ses racines égales, ce qui caractérise le cas additif.

c) \quad Lorsque $p=2$, $\Delta=0$, $c_4\ne 0$, la partie principale de (4.1) en $(0,0)$ est $y^2+xy+a_2x^2=0$, ce qui correspond à un point double à tangentes distinctes.

d) \quad Lorsque $p=2$, $\Delta=0$, $c_{4\unsure{.}}=0$, la partie principale de (6.2) en $(0,0)$ est $y^2+a_4x+a_6=0$, ce qui correspond à une singularité cuspidale.

Soit alors $E$ une courbe de genre un sur une base $S$, sans fibre de type additif. Puisque $\Delta$ et $\unsure{\mathbb{\text{\makebox[0pt][l]{\kern.45pt$c$}$c$}}_4|c_4}$ ne sont jamais simultanément nuls, et sont respectivement des sections de $\omega^{\otimes 12}$ \tikz[baseline=(et.base)]{\node[inner sep=0pt] (et) {et};\handwritten{\draw[thin,gray,rotate around={-27:(et.center)}] (et.center) ellipse (0.80em and 1.16em);}} $\omega^{\otimes 4}$, on définit une section $j$ de
$\tikz[baseline=(P.base)]{\node[inner sep=0pt] (P) {$P_1$};\handwritten{\draw[thin,gray,rotate around={-16:(P.center)}] (P.center) ellipse (0.90em and 1.24em);}\node[inner sep=0pt,below=0.95em] (corr) at (P.south) {\inserted{\handwritten{$P^1$}}};\handwritten{\draw[thin,gray,rotate around={-21:(corr.center)}] (corr.center) ellipse (1.08em and 1.29em);}}/S$
par la formule%
% archive sheet number: 25
\scanpage{27}%
\mspage{12}%
\margin{\handwritten{$\begin{aligned}j&=0 &&\Longleftrightarrow c_4=0\\j&=12^3 &&\Longleftrightarrow c_6=0\end{aligned}$}}%
\[
j=c_4^3/\Delta=1728+c_6^2/\Delta \, . \tag*{(6.2)}
\]

\uline{Proposition} 6.3. \quad \uline{Soient} $E$ \uline{et} $E'$ \uline{deux courbes de genre un sur} $S$, \uline{sans fibre additive} .

(I) $\text{\uline{Isom}}_{\handwritten{S}}(E,E')$ \uline{est représentable, fini et} \tikz[baseline=(nrone.base)]{\node[inner sep=0pt] (nrone) {\uline{non ramifié}};\handwritten{\draw[thin,gray] (nrone.north east) .. controls ([xshift=-1em,yshift=0.8em]nrone.north east) and ([xshift=-0.6em,yshift=0.7em]nrone.north west) .. ([xshift=-0.2em]nrone.west) .. controls ([yshift=-0.85em]nrone.south west) and ([xshift=0.5em,yshift=-0.85em]nrone.south east) .. (nrone.north east);}\node[anchor=south west,inner sep=0pt] at ([xshift=-1.8em,yshift=0.48em]nrone.north east) {\inserted{\handwritten{\unsure{net|met}}}};} \uline{sur} $S$\tikz[baseline=(stop.base)]{\node[inner sep=0pt] (stop) {.};\handwritten{\draw[thin] ([yshift=-0.08em]stop.south) .. controls ([xshift=0.04em,yshift=-0.19em]stop.south) .. ([xshift=-0.15em,yshift=-0.35em]stop.south);\draw[thin,gray,rotate around={-18:(stop.center)}] (stop.center) ellipse (0.56em and 1.34em);}}

(II) \uline{si} $j_E=j_F$, \uline{la projection de} $\text{\uline{I\mbox{\unsure{$\mathbb{\text{\makebox[0pt][l]{\kern.45pt$s$}$s$}}$|s}}om}}_{\handwritten{S}}(E,E')$ \uline{sur} $S$ \uline{est surjective}\tikz[baseline=(stop.base)]{\node[inner sep=0pt] (stop) {.};\handwritten{\draw[thin] ([yshift=-0.08em]stop.south) .. controls ([xshift=0.04em,yshift=-0.19em]stop.south) .. ([xshift=-0.15em,yshift=-0.35em]stop.south);\draw[thin,gray,rotate around={-18:(stop.center)}] (stop.center) ellipse (0.56em and 1.34em);}}

(III) \uline{si} $j_E=j_F$ \uline{ne prend pas les valeurs} $0$ \uline{et} $\underset{\inserted{\handwritten{\substack{\Vert\\12^3}}}}{\tikz[baseline=(val.base)]{\node[inner sep=0pt] (val) {$1728$};\handwritten{\draw[thin] (val.south west) .. controls ([yshift=-0.15em]val.south west) and ([yshift=-0.15em]val.south east) .. (val.south east);}}}$, $\text{\uline{Isom}}_{\handwritten{S}}(E,E')$ \uline{est un revêtement étale de rang} $2$ \uline{de} $S$ .

Donnons-nous $E$ et $E'$ par des équations (1.3) et $(1.3)'$. Se donner un isomorphisme entre $E$ et $E'$ revient à se donner $u$, $r$, $s$ et $t$, vérifiant les équations (1.6), donc aussi (1.7) et (1.8), $u$ devant de plus être inversible. Par hypothèse, localement sur $S$, $c_4$ ou $\Delta$ est inversible (6.1), et de même $c'_4$ ou $\Delta'$ l'est. Localement sur $S$, les équations (1.8) sont donc des équations de dépendance intégrale pour $u$ et $u^{-1}$, les équations $(1.7)_6$ ou $(1.7)_8$ des équations de dépendance intégrale pour $r$ et enfin $(1.6)_6$ (resp. $(1.6)_2$) est une équation de dépendance intégrale pour $t$ (resp. $s$). Ceci prouve que $\text{\uline{Isom}}(E,E')$ est représentable par un schéma affine fini sur $S$ ; il est même \tikz[baseline=(nrtwo.base)]{\node[inner sep=0pt] (nrtwo) {non ramifié};\handwritten{\draw[thin,gray] (nrtwo.north east) .. controls ([xshift=-1em,yshift=0.8em]nrtwo.north east) and ([xshift=-0.6em,yshift=0.7em]nrtwo.north west) .. ([xshift=-0.2em]nrtwo.west) .. controls ([yshift=-0.85em]nrtwo.south west) and ([xshift=0.5em,yshift=-0.85em]nrtwo.south east) .. (nrtwo.north east);}\node[anchor=south west,inner sep=0pt] at ([xshift=-1.8em,yshift=0.48em]nrtwo.north east) {\inserted{\handwritten{net}}};} car les courbes considérées n'ont pas d'automorphismes infinitésimaux (respectant l'origine).

Lorsque $2$ est inversible sur $S$, l'assertion (III) résulte de l'énoncé plus précis.

\uline{Proposition} 6.4. \quad \uline{Soient} $E$ \uline{et} $E'$ \uline{deux courbes de genre un sur} $S$ ,%
% archive sheet number: 26
\scanpage{28}%
\mspage{13}\space
\uline{sans fibre additive}, \uline{munies de formes différentielles invariantes} $\pi$ \uline{et} $\pi'$, \uline{et de même invariant modulaire}. \uline{On suppose que} $2$, $c_4$ \uline{et} $c_6$ \uline{sont inversibles}, \uline{i.e.} \uline{que} $2^{-1}$, $1/j$ et $1/\inserted{\handwritten{(}}j-1728\inserted{\handwritten{)}}$ \uline{sont des sections de} \tikz[baseline=(ess.base)]{\node[inner sep=0pt] (ess) {$S$};\handwritten{\draw[thin,gray,rotate around={-20:(ess.center)}] (ess.center) ellipse (0.52em and 1.02em);}}\reversemarginpar\margin{\inserted{\handwritten{\tikz[baseline=(oh.base)]{\node[inner sep=0pt] (oh) {$\underline{O}_S$};\handwritten{\draw[thin,gray,rotate around={-15:(oh.center)}] (oh.center) ellipse (1.13em and 1.55em);}}}}}\normalmarginpar. \uline{Alors}, \uline{le schéma} $\text{\uline{Isom}}(E,F)$ \uline{s'identifie au schéma défini par l'équation}
\[
\lambda^2=c_4\,c'_6/c_6\,c'_4 \, .
\]

En vertu de (3.6), se donner $(E,\pi)$ revient à se donner $a_2$, $a_4$ et $a_6$ vérifiant $a_2\,a_4-9\,a_6=0$, et un isomorphisme de $E$ sur $E'$ est décrit par $\lambda$ tel que $(E\,\pi)$ soit isomorphe à $(E',\lambda\,\pi')$, i.e. par $\lambda$ vérifiant
\[
a'_2=\lambda^2\,a_2 \qquad a'_4=\lambda^4\,a_4 \qquad a'_6=\lambda^6\,a_6 \, . \tag*{(6.4.1)}
\]

Par hypothèse, $j=j'$, de sorte que $(c_4^3,c_6^2,\Delta)$ est proportionnel à $(c_4^{\raisebox{-.38ex}{$\scriptstyle\prime$}\,3},c_6^{\raisebox{-.38ex}{$\scriptstyle\prime$}\,2},\Delta')$, et que $\mu=c_4\,c'_6/\raisebox{-.52ex}{$c_6\,c'_4$}$ vérifie
\[
c'_4=\mu^2\,c_4 \qquad c'_6=\mu^3\,c_6 \qquad \Delta'=\mu^6\,\Delta \, . \tag*{(6.4.2)}
\]

En vertu de (3.4), on a alors
\[
{\unsure{\mathbb{\text{\makebox[0pt][l]{\kern.45pt$a$}$a$}}|a}}'_2=\mu\,a_2 \qquad a'_4=\mu^2\,a_4 \qquad a'_6=\mu^3\,a_6 \, .
\]

De plus, d'après (3.4) encore, localement, soit $a_2$, soit $a_4$ et $a_6$ sont inversibles, de sorte que (6.4.1) équivaut à $\lambda^2=\mu$.

Pour prouver (III) en général, on peut se ramener à supposer $S$ noethérien, et supposer soit $2$ inversible (cas traité en (6.4)), soit%
% archive sheet number: 27
\scanpage{29}%
\mspage{14}\space
$2$ nilpotent, ce qui permet de mettre la courbe sous la forme (4.1)
\[
y^2+xy=x^3+a_2\,x+a_6
\]
et un isomorphisme d'une telle courbe sur une autre est décrit par $s$ vérifiant (4.3) pour $u=1+2s$. Lorsque $j=j'$, la première de ces équations
\[
(1+2s)^2\,a'_2=a_2-s-s^2
\]
définit un revêtement étale de degré deux, et elle implique la seconde ; pour le voir, on se ramène à supposer $a_2=a'_2$ et on remarque que d'après (4.4)
\[
1/j=a_6+2\,\unsure{\mathbb{R}|R}(a_2,a_6)
\]
pour une fraction rationnelle $R$ .

Pour vérifier (II) dans les cas non couverts par (III), on peut supposer que $S$ est le spectre d'un corps algébriquement clos et que $c_4=0$ ou $c_6=0$ \unsure{,|;} i.e. que $j=0$ ou $j=1728$. Soit $p$ la caractéristique du corps,

a) $p\ne 2,3$ : on met la courbe sous la forme (2.1) ; si $c_4=0$ (resp. $c_6=0$) on aura $a_4=0$ (resp. $a_6=0$) et on peut par homothéti\unsure{c|e} transformer la courbe en la courbe type
\[
y^2=x^3+1 \qquad (\text{si }c_4=0) \tag*{(6.5)}
\]
\[
y^2=x^3-x \qquad (\text{si }c_6=0) \tag*{(6.6)}
\]

b) $\unsure{\mathbb{\text{\makebox[0pt][l]{\kern.45pt$p$}$p$}}|p}=3$ : on met la courbe sour la forme (6.1) ; résolvant (6.2) pour $a'_4=1$ et $a'_6=0$, on met la courbe sous la forme type%
% archive sheet number: 28
\scanpage{30}%
\mspage{15}%
\[
y^2=x^3-x \qquad (\text{si }c_4=c_6=0) \tag*{(6.7)}
\]

c) $p=2$ : ici, on met la courbe sour la forme (6.3) et on résoud (6.4) avec $a_3\unsure{\cdot}=1$, $a_4=a_6=0$ pour obtenir la forme type
\[
y^2+y=x^3 \qquad (\text{si }c_4=c_6=0) \, . \tag*{(6.8)}
\]

Pour compléter le tableau, reste à décrire les automorphismes des courbes elliptiques.

\uline{Proposition} 6.9. \quad \uline{Soit} $E$ \uline{une courbe de genre un sur une base} $S$, \uline{sans fibre additive}.

(I) \quad \uline{Si} $c_4$ et $c_6$ \uline{sont inversibles,} \uline{les seuls automorphismes}\margin{\tikz[baseline=(rem.base)]{\node[align=left,inner sep=2pt] (rem) {\handwritten{\unsure{faux} si $S$}\\\handwritten{\unsure{pas} connexe}};\handwritten{\draw[thin,gray,rounded corners=4pt] (rem.north east) -- (rem.south east) -- (rem.south west) -- (rem.north west) -- ([xshift=-0.4em]rem.north east);\draw[thin] ([xshift=0.35em,yshift=0.5em]rem.north west) -- ([xshift=0.2em]rem.west);}}} \uline{de} $E$ \uline{sont l'identité et la symétrie}.

(II) \quad \uline{Si} $2$ et $3$ \uline{sont inversibles}, \uline{l'application canonique}
\diagram{\[
\vcenter{\hbox{\begin{tikzcd}[ampersand replacement=\&,column sep=1.9em]
 \text{\uline{Aut}}(E) \handwritten{\arrow[r]} \& G_m
\end{tikzcd}}}\tag*{(6.6.)}
\]}
\uline{donnée par l'action sur les différentielles invariantes}, \uline{est une immersion fermée}.

(III) \uline{Si} $2$ \uline{et} $3$ \uline{sont inversibles}, \uline{et} $c_4=0$ (resp. $c_6=0$) \uline{l'image de} (6.6) \uline{est} $\mu_6$ (resp. $\mu_4$).

\uline{Supposons maintenant que} $S$ \uline{soit le spectre d'u\unsure{r|n} corps}\tikz[baseline=(stop.base)]{\node[inner sep=0pt] (stop) {.};\handwritten{\draw[thin,gray,rotate around={-18:(stop.center)}] (stop.center) ellipse (0.57em and 1.03em);}}

(IV) \quad \uline{Si} $3=c_4=c_6=0$, $\mathrm{Aut}(E)$ \uline{est géométriquement \unsure{isomorphe} au produit semi-direct} (\uline{non commutatif}) \uline{de} $\mathbb{Z}/(4)$ \uline{par} $\mathbb{Z}/(3)$, \uline{et l'image de} (6.6) \uline{est} $\mu_4$ .

(V) \quad \uline{Si} $2=c_4=c_6=0$, $\mathrm{Aut}(E)$ \uline{est géométriquement d'ordre} \tikz[baseline=(order.base)]{\node[inner sep=0pt] (order) {$24$,};\handwritten{\draw[thin,gray,rotate around={-38:(order.center)}] ([xshift=.37em,yshift=2.08em]order.center) arc[start angle=72,end angle=365,x radius=1.19em,y radius=2.2em];}}\margin{\handwritten{\tikz[baseline=-0.3ex]{\draw[thin] (0,0.7em) -- (0.35em,-0.05em);} structure ?}} \uline{et l'image de} (6.6) \uline{est} $\mu_3$.
% archive sheet number: 29

\scanpage{31}
\mspage{16}

\uline{Preuve} : (I) \quad D'après (6.3), il n'y a en effet pas plus de deux automorphismes.

(II) et (III) \quad Lorsque les courbes sont mises sous la forme (2.1), il résulte de (2.5) que les automorphismes sont décrits parleur action sur $\pi$ et que $\lambda\in G_m$ défini un automorphisme si et seulement si
\[
c_4=\lambda^4\,c_4 \qquad\text{et}\qquad c_6=\lambda^6\,c_6 \, ,
\]
d'où les assertions.

(IV) \quad Si la courbe est mise sous forme (6.1), les automorphismes sont donnés par (1.2) : $x'\handwritten{\mathrel{\tikz[baseline=-.32ex]{\draw[->,thin] (0,0) .. controls (.19em,.39ex) and (.28em,.6ex) .. (.57em,.18ex) .. controls (.78em,-.23ex) and (1.02em,-.08ex) .. (1.28em,0) -- (2.28em,.04ex);}}}x$, pour $s=t=0$ et
\[
\handwritten{\left\{}\begin{aligned}
&u^4=1\\
&r^3+ra_4+a_6(1-u^2)=0
\end{aligned}\right.
\]
et la loi de composition par
\[
(u',r')\quad (u,r)=(u\,u',u^{\raisebox{-.38ex}{$\scriptstyle\prime$}\,2}r+r')
\]

Lorsque la courbe est mise sous forme (6.7), ces équations se réduisent à $u\in\mu_4$ et $r\in F_3$ .

(V) \quad Si la courbe est mise sous forme (6.3), les automorphismes sont donnés par (1.2) : $x'\handwritten{\mathrel{\tikz[baseline=-.32ex]{\draw[->,thin] (0,0) .. controls (.14em,.37ex) and (.25em,.52ex) .. (.44em,.13ex) .. controls (.62em,-.25ex) and (.86em,-.2ex) .. (1.09em,.08ex) .. controls (1.3em,.31ex) and (1.52em,.23ex) .. (1.67em,.20ex) -- (2.36em,.16ex);}}}x$, pour $r=s^2$ et
\[
\handwritten{\left\{}\begin{aligned}
&u^3=1\\
&s^4+sa_3+a_4(1-u)=0\\
&t^2+ta_3+s^6+s^2\,a_4=0
\end{aligned}\right.
\]%
% archive sheet number: 30
\scanpage{32}%
\mspage{17}%
et la loi de composition par
\[
(u',s',t')\mathbin{\handwritten{\unsure{o|\circ}}}(u,s,t)
=(u\,u',u's+t',u^{\raisebox{-.38ex}{$\scriptstyle\prime$}\,3}t+t'+u^{\raisebox{-.38ex}{$\scriptstyle\prime$}\,2}rs')\, .
\]

Lorsque la courbe est mise sous forme (6.8), ces équations se réduisent à
\[
\begin{aligned}
u\in\mu_3 \qquad s\in\mathbb{F}_4 \qquad t^2+t=s^3
 &\quad ({}=0\ \text{ou}\ 1\ \text{selon que}\\
 &\qquad s=0\ \text{ou}\ s\ne 0) \, .
\end{aligned}
\]

7. \quad \uline{Anneau des formes modulaires entières}.

On appelle forme modulaire entière de poids $n$ une loi qui à chaque courbe de genre un sur une base $S$ associe une section de $\omega^{\otimes n}$, et ce de façon compatible au changement de base. Il reviendrait\margin{\handwritten{\tikz[baseline=-0.2ex]{\draw[thin,gray] (0,-0.35em) -- (0.48em,1.30em);} \uuline{pourquoi} ?}} au même de se limiter aux courbes sans fibres additives.

\uline{Proposition} 7.1. \quad \uline{L'anneau des formes modulaires entières est engendré sur} $\mathbb{Z}$ \uline{par} $c_4$, $c_6$ \uline{et} $\Delta$ \uline{soumis à la seule relation}
\[
c_4^3-c_6^2=1728\,\Delta \, . \tag*{(7.2.)}
\]
\margin{\tikz[baseline=(note.north)]{\node[inner sep=2pt,text width=9.8em,align=left] (note) {\handwritten{N.B. Sur le site modulaire, les diviseurs de $c_4$, $c_6$, $\Delta$ sont à \unsure{multiplicité|multiplicités} $1$ ; \unsure{Celui|celui} de $\Delta$ est lisse sur $\mathbb{Z}$, ceux de $c_4$, $c_6$ le sont sauf au dessus des points $2$ et $3$.}};\handwritten{\draw[thin,rounded corners=1pt] (note.north west) -- (note.south west) -- (note.south east);}}}

Appliquant la définition à la courbe (1.3), on voit qu'une forme modulaire entière de poids $n$ s'exprime comme poly\tikz[baseline=(nome.base)]{\node[inner sep=0pt] (nome) {nôme};\handwritten{\draw[thin,gray,rotate around={-25:(nome.center)}] (nome.center) ellipse (1.04em and 1.65em);}} de poids $n$ en les $a_i$. Pour $2$ et $3$ inversibles, il résulte de 2.5 que ces formes sont \tikz[baseline=(des.base)]{\node[inner sep=0pt] (des) {des};\handwritten{\draw[thin,gray,rotate around={-15:(des.center)}] (des.center) ellipse (1em and 1.02em);}} poly\tikz[baseline=(noms.base)]{\node[inner sep=0pt] (noms) {nôm};\handwritten{\draw[thin,gray,rotate around={-25:(noms.center)}] (noms.center) ellipse (0.84em and 1.47em);}}es en $c_4$ et $c_6$, et en particulier que $c_4$, $c_6$ et $\Delta$ ne satisfont pas d'autres relations que (7.2). Reste à prouver que si un poly\tikz[baseline=(nomt.base)]{\node[inner sep=0pt] (nomt) {nôme};\handwritten{\draw[thin,gray,rotate around={-25:(nomt.center)}] (nomt.center) ellipse (0.85em and 1.53em);}} en $c_4$, $c_6$ et $\Delta$ est identiquement nul en caractéristique $2$ (resp. $3$), alors il est divisible par $2$ (resp.$3$) dans $\mathbb{Z}[c_4,c_6,\Delta]/\handwritten{\bigl(}c_4^3-c_6^2-1728\,\Delta\handwritten{\bigr)}$, ce qui signifie%
% archive sheet number: 31
\scanpage{33}%
\mspage{18}\space
qu'il est divisible par $c_4^3-c_6^2$ dans $\mathbb{Z}/(2)[c_4,c_6,\Delta]$ (resp. $\mathbb{Z}/(3)[c_4,c_6,\Delta]$). Cela résulte de la structur\unsure{e|c} de l'anneau des formes modulaires de caractéristique $2$ (resp. $3$).

\uline{Proposition} 7.2. \quad (I) \uline{L'anneau des formes modulaires de caractéristique} $2$ \uline{est} $\mathbb{Z}/(2)[\tikz[baseline=(coeffa.base)]{\node[inner sep=0pt] (coeffa) {$a_1$};\handwritten{\node[inner sep=0pt,anchor=west,font=\scriptsize] (notea) at ([xshift=1.35em,yshift=-1.05em]coeffa.south west) {le diviseur de $a_1$ est lisse sur $\mathbb{F}_2$};\draw[thin] (coeffa.south west) to[out=-86,in=192] ([xshift=-.18em]notea.west);}\pgfresetboundingbox\path[use as bounding box] ([yshift=-1.5em]coeffa.south west) rectangle (coeffa.north east);},\Delta]$ , \uline{et} $c_4=a_1^4$ , $c_6=a_1^6$ .

\hspace*{3em}(II) \uline{L'anneau des formes modulaires de caractéristique} $3$ \uline{est} $\mathbb{Z}/(3)[\tikz[baseline=(coeffb.base)]{\node[inner sep=0pt] (coeffb) {$b_2$};\handwritten{\node[inner sep=0pt,anchor=west,font=\scriptsize] (noteb) at ([xshift=1.35em,yshift=-1.05em]coeffb.south west) {le diviseur de $b_2$ est lisse sur $\mathbb{F}_2$};\draw[thin] (coeffb.south west) to[out=-86,in=192] ([xshift=-.18em]noteb.west);}\pgfresetboundingbox\path[use as bounding box] ([yshift=-1.5em]coeffb.south west) rectangle (coeffb.north east);},\Delta]$ , et $c_4=b_2^2$ , $c_6=-b_2^3$ .

Tout d'abord, il résulte de $(1.6)_1$ et $(1.7)_2$ que $a_1$ (resp. $b_2$) est une forme modulaire de caractéristique $2$ (resp. $3$) .

En caractéristique $2$ , cherchons tout d'abord les formes modulaires relatives aux seules courbes telles que $a_1$ soit inversible. Ces courbes se mettent sous forme (cf. (4.\tikz[baseline=(un.base)]{\node[inner sep=0pt] (un) {1};\handwritten{\draw[thin,gray,rotate around={-12:(un.center)}] (un.center) ellipse (0.36em and 1.12em);}}))
\[
 y^2+a_1xy=x^3+a_2x^2+a_6 \qquad \pi=\frac{dx}{a_1x}
\]
et, compte tenu de (4.3), on voit que ces formes s'écrivent
\[
\begin{aligned}
 &a_1^{-n}\,P(a_1,a_6) \qquad \text{pour }P\text{ isobare, ou encore}\\[.65em]
 &a_1^{-m}\,Q(a_1,\Delta)\quad\handwritten{,}
\end{aligned}
\]
et cela garde un sens pour $a_1=0$ , $\Delta\ne0$ si et seulement si on avait en fait affaire à un polyn\tikz[baseline=(oc.base)]{\node[inner sep=0pt] (oc) {ô};\handwritten{\draw[thin,gray,rotate around={-14:(oc.center)}] (oc.center) ellipse (0.56em and 1.1em);}}me en $a_1$ et $\Delta$ .

En caractéristique $3$ , raisonnant de même d'abord sur les courbes telles que $c_4$ soit inversible, et remarquant que $d_6=0$ équivaut alors à $a_4=0$ dans (3.1), on voit que toute forme modu%
% archive sheet number: 32
\scanpage{34}%
\mspage{19}%
laire s'écrit
\[
 b_2^{-n}\,P(b_2,\Delta) \qquad \text{pour }P\text{ isobare}\quad ,
\]
et que cette forme modulaire est définie aussi pour $j=0$ si et seulement si elle s'écrit comme polyn\tikz[baseline=(oend.base)]{\node[inner sep=0pt] (oend) {ô};\handwritten{\draw[thin,gray,rotate around={-20:(oend.center)}] (oend.center) ellipse (0.64em and 1.15em);}}me en $b_2$ et $\Delta$ .

\handwritten{Interpr\unsure{é|e}tation de $a_1$, $b_2$ comme forme de Hasse-Witt ?}

\handwritten{\unsure{Formes modulaires|Forme modulaire} \unsure{sur|en} un corps quelconque ? Sur un ann. de val. discrète\ldots}

\handwritten{Hasse-Witt sur un corps de \unsure{car|car.} $p>0$ quelconque ?}

\noindent\handwritten{\tikz{\draw[thin] (0,0) -- (3em,.06em);}}

\[
\handwritten{\left\{\begin{aligned}
 &\text{\unsure{Relèvement|Relévement} canonique de Serre par formules universelles}\,\big|\\
 &\text{Loi de groupe formel sur }{\unsure{\mathcal{P}|P}}^{\infty}_{\unsure{M|X}/S}\ \text{: formules explicites\ldots}
\end{aligned}\right.}
\]
% archive sheet number: 33

